% Appendices are numbered A, B, ... and headed "APPENDIX A".
\titleformat{\chapter}[display]{\centering\normalfont\bfseries}{\Large APPENDIX \thechapter}{6pt}{\LARGE}
\renewcommand{\chaptermark}[1]{\markboth{Appendix \thechapter}{}}% short: the long titles collide with the left header

\chapter{Folder Map}\label{app:foldermap}\label{app:A}

This appendix tells a reader where to find, in the project folder \code{tieuluan/}, every part of the work that
this essay describes. All paths are relative to that folder. The tables are generated when the essay is built
(\code{report/build.sh}), so they always match the files that exist; a figure marked ``missing'' has not been
exported yet.

\section{Chapters and their sources}\label{sec:A.1}

\begin{longtable}{@{}P{5.0cm}P{6.2cm}P{3.8cm}@{}}
\caption{Chapters of the essay and their LaTeX sources}\label{tab:mapchapters}\\
\toprule
\textbf{Part of the essay} & \textbf{LaTeX source} & \textbf{Main inputs}\\
\midrule
\endfirsthead
\toprule
\textbf{Part of the essay} & \textbf{LaTeX source} & \textbf{Main inputs}\\
\midrule
\endhead
\small Cover, contents, lists & \small \code{report/main.tex}, \code{report/chapters/cover.tex}, \code{report/members.tex} & \small Group information\\
\small Abbreviations, Introduction & \small \code{report/chapters/front.tex} & \small \code{tools/model.py}\\
\small Chapter 1 Requirement Determination & \small \code{report/chapters/ch1_requirements.tex} & \small \code{tools/model.py}, \code{figures/vp/}\\
\small Chapter 2 Analysis & \small \code{report/chapters/ch2_analysis.tex} & \small \code{tools/model.py}, \code{figures/vp/}\\
\small Chapter 3 AI in E-commerce & \small \code{report/chapters/ch3_ai.tex} & \small \code{build/eval/}, \code{figures/}\\
\small Chapter 4 Design & \small \code{report/chapters/ch4_design.tex} & \small \code{tools/model.py}, \code{demo/db/}, \code{demo/generated/}\\
\small Chapter 5 Code & \small \code{report/chapters/ch5_code.tex} & \small \code{demo/}, \code{figures/demo/}\\
\small Conclusion, References & \small \code{report/chapters/conclusion.tex}, \code{report/chapters/references.tex} & \small --\\
\small Appendix A Folder Map & \small \code{report/chapters/appendix.tex} & \small \code{report/gen_tables.py}\\
\small Preamble, build & \small \code{report/preamble.tex}, \code{report/build.sh}, \code{report/gen_tables.py} & \small --\\
\bottomrule
\end{longtable}

\section{Figures, code excerpts, folders and generated tables}\label{sec:A.2}

Table~\ref{tab:mapfigs} lists every figure with the file it is taken from; diagrams are exported from the Visual
Paradigm project in \code{vp/}, screenshots are taken from the running demo. Table~\ref{tab:maplists} gives the real
source file of every code excerpt, Table~\ref{tab:mapfolders} the main folders of the project, and
Table~\ref{tab:mapgen} the tables that are generated from the shared specification and the code.

\begin{longtable}{@{}P{1.7cm}P{10.4cm}P{2.6cm}@{}}
\caption{Figures of the essay and their files}\label{tab:mapfigs}\\
\toprule
\footnotesize \textbf{Chapter} & \footnotesize \textbf{File} & \footnotesize \textbf{Status}\\
\midrule
\endfirsthead
\toprule
\footnotesize \textbf{Chapter} & \footnotesize \textbf{File} & \footnotesize \textbf{Status}\\
\midrule
\endhead
\bottomrule
\endlastfoot
\footnotesize 2 & \footnotesize \code{figures/vp/c2_uc_search.png} & \footnotesize $\checkmark$\\
\footnotesize 2 & \footnotesize \code{figures/vp/c2_uc_purchase.png} & \footnotesize $\checkmark$\\
\footnotesize 2 & \footnotesize \code{figures/vp/c2_uc_assistant_admin.png} & \footnotesize $\checkmark$\\
\footnotesize 2 & \footnotesize \code{figures/vp/c2_activity_search.png} & \footnotesize $\checkmark$\\
\footnotesize 2 & \footnotesize \code{figures/vp/c2_bce_search.png} & \footnotesize $\checkmark$\\
\footnotesize 2 & \footnotesize \code{figures/vp/c2_seq_voice.png} & \footnotesize $\checkmark$\\
\footnotesize 2 & \footnotesize \code{figures/vp/c2_activity_image.png} & \footnotesize $\checkmark$\\
\footnotesize 2 & \footnotesize \code{figures/vp/c2_activity_checkout.png} & \footnotesize $\checkmark$\\
\footnotesize 2 & \footnotesize \code{figures/vp/c2_bce_checkout.png} & \footnotesize $\checkmark$\\
\footnotesize 2 & \footnotesize \code{figures/vp/c2_activity_cancel.png} & \footnotesize $\checkmark$\\
\footnotesize 2 & \footnotesize \code{figures/vp/c2_bce_cancel.png} & \footnotesize $\checkmark$\\
\footnotesize 2 & \footnotesize \code{figures/vp/c2_state_order.png} & \footnotesize $\checkmark$\\
\footnotesize 2 & \footnotesize \code{figures/vp/c2_activity_admin.png} & \footnotesize $\checkmark$\\
\footnotesize 2 & \footnotesize \code{figures/vp/c2_domain.png} & \footnotesize $\checkmark$\\
\footnotesize 3 & \footnotesize \code{figures/vp/c3_ai_architecture.png} & \footnotesize $\checkmark$\\
\footnotesize 3 & \footnotesize \code{figures/demo/search_voice.png} & \footnotesize $\checkmark$\\
\footnotesize 3 & \footnotesize \code{figures/demo/search_image.png} & \footnotesize $\checkmark$\\
\footnotesize 3 & \footnotesize \code{figures/vp/c3_seq_image.png} & \footnotesize $\checkmark$\\
\footnotesize 3 & \footnotesize \code{figures/demo/room_preview.png} & \footnotesize $\checkmark$\\
\footnotesize 3 & \footnotesize \code{figures/vp/c3_seq_assistant.png} & \footnotesize $\checkmark$\\
\footnotesize 3 & \footnotesize \code{figures/demo/assistant.png} & \footnotesize $\checkmark$\\
\footnotesize 3 & \footnotesize \code{figures/demo/ai_fallback.png} & \footnotesize $\checkmark$\\
\footnotesize 3 & \footnotesize \code{figures/demo/search_debug.png} & \footnotesize $\checkmark$\\
\footnotesize 4 & \footnotesize \code{figures/demo/db_tables.png} & \footnotesize $\checkmark$\\
\footnotesize 4 & \footnotesize \code{figures/demo/db_rows.png} & \footnotesize $\checkmark$\\
\footnotesize 4 & \footnotesize \code{figures/vp/c44_packages.png} & \footnotesize $\checkmark$\\
\footnotesize 4 & \footnotesize \code{figures/vp/c44_checkout_seq.png} & \footnotesize $\checkmark$\\
\footnotesize 4 & \footnotesize \code{figures/vp/c44_seq_search.png} & \footnotesize $\checkmark$\\
\footnotesize 4 & \footnotesize \code{figures/vp/c44_seq_cancel.png} & \footnotesize $\checkmark$\\
\footnotesize 5 & \footnotesize \code{figures/vp/c5_deployment.png} & \footnotesize $\checkmark$\\
\footnotesize 5 & \footnotesize \code{figures/demo/home.png} & \footnotesize $\checkmark$\\
\footnotesize 5 & \footnotesize \code{figures/demo/search_text.png} & \footnotesize $\checkmark$\\
\footnotesize 5 & \footnotesize \code{figures/demo/product_detail.png} & \footnotesize $\checkmark$\\
\footnotesize 5 & \footnotesize \code{figures/demo/cart.png} & \footnotesize $\checkmark$\\
\footnotesize 5 & \footnotesize \code{figures/demo/checkout.png} & \footnotesize $\checkmark$\\
\footnotesize 5 & \footnotesize \code{figures/demo/orders.png} & \footnotesize $\checkmark$\\
\footnotesize 5 & \footnotesize \code{figures/demo/order_forbidden.png} & \footnotesize $\checkmark$\\
\footnotesize 5 & \footnotesize \code{figures/demo/order_cancel.png} & \footnotesize $\checkmark$\\
\footnotesize 5 & \footnotesize \code{figures/demo/admin_products.png} & \footnotesize $\checkmark$\\
\footnotesize 5 & \footnotesize \code{figures/demo/admin_orders.png} & \footnotesize $\checkmark$\\
\end{longtable}
\begin{longtable}{@{}P{1.7cm}P{5.4cm}P{7.9cm}@{}}
\caption{Code excerpts of the essay and their real source files}\label{tab:maplists}\\
\toprule
\footnotesize \textbf{Chapter} & \footnotesize \textbf{Excerpt} & \footnotesize \textbf{Source file}\\
\midrule
\endfirsthead
\toprule
\footnotesize \textbf{Chapter} & \footnotesize \textbf{Excerpt} & \footnotesize \textbf{Source file}\\
\midrule
\endhead
\bottomrule
\endlastfoot
\footnotesize 4 & \footnotesize \code{build/excerpts5/CheckoutService.java} & \footnotesize \code{demo/src/main/java/com/ptit/ecom/application/service/CheckoutService.java}\\
\footnotesize 5 & \footnotesize \code{build/excerpts5/EcomApplication.java} & \footnotesize \code{demo/src/main/java/com/ptit/ecom/presentation/web/EcomApplication.java}\\
\footnotesize 5 & \footnotesize \code{build/excerpts5/QueryService.java} & \footnotesize \code{demo/src/main/java/com/ptit/ecom/application/search/QueryService.java}\\
\footnotesize 5 & \footnotesize \code{build/excerpts5/SearchService.java} & \footnotesize \code{demo/src/main/java/com/ptit/ecom/application/search/SearchService.java}\\
\footnotesize 5 & \footnotesize \code{build/excerpts5/VectorIndex.java} & \footnotesize \code{demo/src/main/java/com/ptit/ecom/infrastructure/search/VectorIndex.java}\\
\footnotesize 5 & \footnotesize \code{build/excerpts5/RankingService.java} & \footnotesize \code{demo/src/main/java/com/ptit/ecom/application/search/RankingService.java}\\
\footnotesize 5 & \footnotesize \code{build/excerpts5/GemmaClient.java} & \footnotesize \code{demo/src/main/java/com/ptit/ecom/infrastructure/ai/GemmaClient.java}\\
\footnotesize 5 & \footnotesize \code{build/excerpts5/WhisperClient.java} & \footnotesize \code{demo/src/main/java/com/ptit/ecom/infrastructure/ai/WhisperClient.java}\\
\footnotesize 5 & \footnotesize \code{build/excerpts5/AssistantService.java} & \footnotesize \code{demo/src/main/java/com/ptit/ecom/application/service/AssistantService.java}\\
\footnotesize 5 & \footnotesize \code{build/excerpts5/voice_input.js} & \footnotesize \code{demo/web/js/voice-input.js}\\
\end{longtable}
\begin{longtable}{@{}P{4.0cm}P{9.0cm}P{2.0cm}@{}}
\caption{Main folders of the project}\label{tab:mapfolders}\\
\toprule
\footnotesize \textbf{Path} & \footnotesize \textbf{Content} & \footnotesize \textbf{Chapter}\\
\midrule
\endfirsthead
\toprule
\footnotesize \textbf{Path} & \footnotesize \textbf{Content} & \footnotesize \textbf{Chapter}\\
\midrule
\endhead
\bottomrule
\endlastfoot
\footnotesize \code{tools/model.py} & \footnotesize Shared specification: requirements, use cases, classes, relationships, layers, tables & \footnotesize 1--5\\
\footnotesize \code{vp/} & \footnotesize Visual Paradigm project with every diagram & \footnotesize 1--5\\
\footnotesize \code{figures/vp/} & \footnotesize Diagrams exported from Visual Paradigm (PNG, PDF, SVG) & \footnotesize 1--5\\
\footnotesize \code{figures/demo/} & \footnotesize Screenshots of the demo and of the database & \footnotesize 3--5\\
\footnotesize \code{demo/db/} & \footnotesize SQLite schema and seed data generated from tools/model.py & \footnotesize 4\\
\footnotesize \code{demo/generated/src/} & \footnotesize Java skeleton generated from the 4.4 model & \footnotesize 4\\
\footnotesize \code{demo/src/main/java/} & \footnotesize Java source of the application, one folder per package & \footnotesize 4, 5\\
\footnotesize \code{demo/web/} & \footnotesize Browser views (HTML, CSS, JavaScript) and painting images & \footnotesize 5\\
\footnotesize \code{demo/tests/java/} & \footnotesize Automated tests of the demo and their fixtures & \footnotesize 5\\
\footnotesize \code{demo/README.md} & \footnotesize Package map of the demo and how to build, run and test it & \footnotesize 5\\
\footnotesize \code{build/MODEL_CHANGES.md} & \footnotesize Changes to tools/model.py proposed while coding, with reasons & \footnotesize 4, 5\\
\footnotesize \code{demo/data/} & \footnotesize Painting catalogue, image credits, evaluation images & \footnotesize 3, 5\\
\footnotesize \code{demo/voice-server/} & \footnotesize Self-hosted Whisper server for speech-to-text & \footnotesize 3, 5\\
\footnotesize \code{build/eval/} & \footnotesize Evaluation results of search and voice & \footnotesize 3, 5\\
\footnotesize \code{build/excerpts5/} & \footnotesize Code excerpts printed in this essay & \footnotesize 4, 5\\
\footnotesize \code{report/} & \footnotesize LaTeX source of this essay; report/build.sh rebuilds the PDF & \footnotesize --\\
\footnotesize \code{slides/} & \footnotesize Presentation & \footnotesize --\\
\end{longtable}
\begin{longtable}{@{}P{7.0cm}P{8.0cm}@{}}
\caption{Tables generated from the shared specification and the code}\label{tab:mapgen}\\
\toprule
\footnotesize \textbf{Table file} & \footnotesize \textbf{Generated from}\\
\midrule
\endfirsthead
\toprule
\footnotesize \textbf{Table file} & \footnotesize \textbf{Generated from}\\
\midrule
\endhead
\bottomrule
\endlastfoot
\footnotesize \code{report/generated/actors.tex} & \footnotesize \code{tools/model.py}\\
\footnotesize \code{report/generated/analysis_classes.tex} & \footnotesize \code{tools/model.py}\\
\footnotesize \code{report/generated/analysis_rel.tex} & \footnotesize \code{tools/model.py}\\
\footnotesize \code{report/generated/app_classes.tex} & \footnotesize \code{tools/model.py}\\
\footnotesize \code{report/generated/config.tex} & \footnotesize \code{demo/config/app.properties}\\
\footnotesize \code{report/generated/design_attrs.tex} & \footnotesize \code{tools/model.py}\\
\footnotesize \code{report/generated/design_nav.tex} & \footnotesize \code{tools/model.py}\\
\footnotesize \code{report/generated/design_ops.tex} & \footnotesize \code{tools/model.py}\\
\footnotesize \code{report/generated/endpoints.tex} & \footnotesize \code{demo/src/main/java/.../WebServer.java}\\
\footnotesize \code{report/generated/errors.tex} & \footnotesize \code{demo/src/main/java/.../WebServer.java}\\
\footnotesize \code{report/generated/eval_values.tex} & \footnotesize \code{tools/model.py}\\
\footnotesize \code{report/generated/fks.tex} & \footnotesize \code{tools/model.py}\\
\footnotesize \code{report/generated/layers.tex} & \footnotesize \code{tools/model.py}\\
\footnotesize \code{report/generated/listings.tex} & \footnotesize \code{tools/model.py}\\
\footnotesize \code{report/generated/moscow.tex} & \footnotesize \code{tools/model.py}\\
\footnotesize \code{report/generated/nfr_check.tex} & \footnotesize \code{tools/model.py}\\
\footnotesize \code{report/generated/req_fr.tex} & \footnotesize \code{tools/model.py}\\
\footnotesize \code{report/generated/req_nfr.tex} & \footnotesize \code{tools/model.py}\\
\footnotesize \code{report/generated/tables_list.tex} & \footnotesize \code{tools/model.py}\\
\footnotesize \code{report/generated/trace_req_uc.tex} & \footnotesize \code{tools/model.py}\\
\footnotesize \code{report/generated/uc_links.tex} & \footnotesize \code{tools/model.py}\\
\footnotesize \code{report/generated/usecases.tex} & \footnotesize \code{tools/model.py}\\
\footnotesize \code{report/generated/views.tex} & \footnotesize \code{tools/model.py, demo/web/}\\
\footnotesize \code{report/generated/voice_values.tex} & \footnotesize \code{tools/model.py}\\
\end{longtable}


\section{Full file tree of the demo}\label{sec:A.3}

The tree below, written by the demo build (\code{build/code_tree5.txt}), lists every file of \code{demo}; the package
map is explained in \code{demo/README.md}, and the changes to \code{tools/model.py} proposed while coding, with their
reasons, are recorded in \code{build/MODEL_CHANGES.md}.

\txtfile{../build/code_tree5.txt}

\chapter{Design Reference Tables}\label{app:designtables}\label{app:B}

The tables of this appendix complete Chapter~4. They are reference lists rather than explanations: the chapter
shows the diagrams and explains the decisions, and the complete lists are kept here. Like every table that names a
class, a method, a table or a column, they are generated from \code{tools/model.py} when the essay is built.

\section{Classes and relationships of the analysis model}\label{sec:B.1}

Table~\ref{tab:classes41} gives the business meaning and the attributes of the \nAnalysisClasses{} classes of
Figure~\ref{fig:c41_analysis} (Section~4.1), and Table~\ref{tab:rel41} explains each of their relationships.

\begin{longtable}{@{}P{2.6cm}P{5.2cm}P{7.2cm}@{}}
\caption{The analysis classes}\label{tab:classes41}\\
\toprule
\small \textbf{Class} & \small \textbf{Business meaning} & \small \textbf{Attributes}\\
\midrule
\endfirsthead
\toprule
\small \textbf{Class} & \small \textbf{Business meaning} & \small \textbf{Attributes}\\
\midrule
\endhead
\bottomrule
\endlastfoot
\small \code{FullName} & \small Name of a customer, a value object & \small \code{firstName}, \code{middleName}, \code{lastName}\\
\small \code{Customer} & \small A customer with an account & \small \code{customerId}, \code{email}, \code{phone}, \code{dateOfBirth}\\
\small \code{Address} & \small Contact and delivery address & \small \code{houseNumber}, \code{street}, \code{district}, \code{city}\\
\small \code{CustomerVIP} & \small VIP customer with a discount and reward points & \small \code{vipLevel}, \code{discountRate}, \code{rewardPoints}\\
\small \code{CustomerNew} & \small Newly registered customer with a welcome voucher & \small \code{registeredDate}, \code{welcomeVoucher}\\
\small \code{Cart} & \small Current shopping state (temporary) & \small \code{cartId}, \code{createdDate}, \code{updatedDate}\\
\small \code{CartItem} & \small A product and the intended quantity in the cart & \small \code{quantity}, \code{cartPrice}\\
\small \code{Order} & \small A confirmed purchase transaction & \small \code{orderId}, \code{orderDate}, \code{status}, \code{totalAmount}\\
\small \code{OrderItem} & \small A purchased product with its frozen transaction price & \small \code{quantity}, \code{unitPrice}\\
\small \code{Shipping} & \small Delivery of one order to one address & \small \code{shippingId}, \code{carrier}, \code{trackingNumber}, \code{shippingFee}, \code{status}, \code{deliveredDate}\\
\small \code{Payment} \emph{(abstract)} & \small One payment attempt (abstract) & \small \code{paymentId}, \code{amount}, \code{paymentDate}, \code{status}\\
\small \code{PayCash} & \small Cash on delivery & \small \code{cashReceived}, \code{changeAmount}\\
\small \code{PayCredit} & \small Card payment through the payment service & \small \code{cardHolder}, \code{bankName}, \code{cardLast4}, \code{authorizationCode}\\
\small \code{Product} \emph{(abstract)} & \small Anything the gallery sells (abstract) & \small \code{productId}, \code{name}, \code{price}, \code{stockQuantity}, \code{imageUrl}, \code{description}\\
\small \code{Artist} & \small Painter of the original work & \small \code{artistId}, \code{name}, \code{nationality}, \code{lifeYears}\\
\small \code{Painting} & \small A fine-art print of one painting & \small \code{style}, \code{subject}, \code{dominantColor}, \code{medium}, \code{yearCreated}, \code{widthCm}, \code{heightCm}\\
\small \code{Frame} & \small A frame sold with a print & \small \code{material}, \code{color}, \code{sizeLabel}\\
\end{longtable}


\begin{longtable}{@{}P{4.3cm}P{2.3cm}P{2.0cm}P{6.4cm}@{}}
\caption{Relationships between the analysis classes}\label{tab:rel41}\\
\toprule
\small \textbf{Relationship (name)} & \small \textbf{Kind} & \small \textbf{Multiplicity} & \small \textbf{Business meaning}\\
\midrule
\endfirsthead
\toprule
\small \textbf{Relationship (name)} & \small \textbf{Kind} & \small \textbf{Multiplicity} & \small \textbf{Business meaning}\\
\midrule
\endhead
\bottomrule
\endlastfoot
\small \code{Customer} -- \code{FullName} (has) & \small Composition & \small 1 -- 1 & \small Each customer has exactly one name; the name does not exist without the customer, hence composition.\\
\small \code{Customer} -- \code{Address} (has) & \small Aggregation & \small 1 -- 1..* & \small A customer keeps at least one delivery address; Shipping also uses it, so only aggregation.\\
\small \code{Customer} -- \code{Cart} (owns) & \small Association & \small 1 -- 0..1 & \small At most one active cart per customer; a new customer may have none.\\
\small \code{Customer} -- \code{Order} (places) & \small Association & \small 1 -- 0..* & \small A customer places many orders; a new customer has none yet.\\
\small \code{Cart} -- \code{CartItem} (contains) & \small Composition & \small 1 -- 0..* & \small A cart line exists only inside its cart; an empty cart is valid.\\
\small \code{Order} -- \code{OrderItem} (contains) & \small Composition & \small 1 -- 1..* & \small An order line is part of its order; an order without lines makes no sense.\\
\small \code{CartItem} -- \code{Product} (refers to) & \small Association & \small 0..* -- 1 & \small Each cart line refers to exactly one product; a product appears in many carts.\\
\small \code{OrderItem} -- \code{Product} (refers to) & \small Association & \small 0..* -- 1 & \small Each order line refers to exactly one product; a product appears in many orders.\\
\small \code{Order} -- \code{Shipping} (is delivered by) & \small Association & \small 1 -- 0..1 & \small A new order may not be shipped yet; once shipped it has exactly one delivery record.\\
\small \code{Shipping} -- \code{Address} (ships to) & \small Association & \small 0..* -- 1 & \small Each delivery goes to exactly one address; an address receives many deliveries.\\
\small \code{Order} -- \code{Payment} (is paid by) & \small Association & \small 1 -- 0..* & \small An order may be unpaid, or have several attempts when one fails.\\
\small \code{Artist} -- \code{Painting} (paints) & \small Association & \small 1 -- 0..* & \small An artist has many paintings in the gallery; each print belongs to exactly one artist.\\
\small \code{Category} -- \code{Product} (classifies) & \small Association & \small 1 -- 0..* & \small Every product is classified in exactly one category; a new category may still be empty.\\
\small \code{CustomerVIP}, \code{CustomerNew} $\rightarrow$ \code{Customer} & \small Generalization & \small  & \small Both kinds of customer reuse the attributes of Customer and add their own.\\
\small \code{PayCash}, \code{PayCredit} $\rightarrow$ \code{Payment} & \small Generalization & \small  & \small Both payment methods share amount, date and status.\\
\small \code{Painting}, \code{Frame} $\rightarrow$ \code{Product} & \small Generalization & \small  & \small Paintings and frames share id, name, price, stock, image and description.\\
\small \code{Cart} $\dashrightarrow$ \code{Order} (checkout) & \small Dependency «create» & \small  & \small Checkout uses the cart content to create a new order; the cart is then emptied.\\
\end{longtable}


\section{Tables and foreign keys of the database}\label{sec:B.2}

Table~\ref{tab:tables} lists the \nTables{} tables of the physical model (Figure~\ref{fig:c42_tables}) with their
primary keys, their foreign keys and the class each table comes from; Table~\ref{tab:fks} gives the name and the
multiplicity of each of the \nFks{} foreign keys (Section~4.2).

\begin{longtable}{@{}P{0.6cm}P{2.8cm}P{3.2cm}P{5.0cm}P{3.4cm}@{}}
\caption{Tables of the database}\label{tab:tables}\\
\toprule
\small \textbf{\#} & \small \textbf{Table} & \small \textbf{Primary key} & \small \textbf{Foreign key $\rightarrow$ parent} & \small \textbf{Source class}\\
\midrule
\endfirsthead
\toprule
\small \textbf{\#} & \small \textbf{Table} & \small \textbf{Primary key} & \small \textbf{Foreign key $\rightarrow$ parent} & \small \textbf{Source class}\\
\midrule
\endhead
\bottomrule
\endlastfoot
\small 1 & \small \code{customer} & \small \code{customer_id} & \small -- & \small Customer + FullName\\
\small 2 & \small \code{customer_vip} & \small \code{customer_id} & \small \code{customer_id} $\rightarrow$ \code{customer} & \small CustomerVIP\\
\small 3 & \small \code{customer_new} & \small \code{customer_id} & \small \code{customer_id} $\rightarrow$ \code{customer} & \small CustomerNew\\
\small 4 & \small \code{address} & \small \code{address_id} & \small \code{customer_id} $\rightarrow$ \code{customer} & \small Address\\
\small 5 & \small \code{category} & \small \code{category_id} & \small -- & \small Category\\
\small 6 & \small \code{product} & \small \code{product_id} & \small \code{category_id} $\rightarrow$ \code{category} & \small Product\\
\small 7 & \small \code{artist} & \small \code{artist_id} & \small -- & \small Artist\\
\small 8 & \small \code{painting} & \small \code{product_id} & \small \code{product_id} $\rightarrow$ \code{product}; \code{artist_id} $\rightarrow$ \code{artist} & \small Painting\\
\small 9 & \small \code{frame} & \small \code{product_id} & \small \code{product_id} $\rightarrow$ \code{product} & \small Frame\\
\small 10 & \small \code{product_feature} & \small \code{product_id} & \small \code{product_id} $\rightarrow$ \code{product} & \small VectorIndex (search)\\
\small 11 & \small \code{cart} & \small \code{cart_id} & \small \code{customer_id} $\rightarrow$ \code{customer} & \small Cart\\
\small 12 & \small \code{cart_item} & \small \code{cart_id}, \code{product_id} & \small \code{cart_id} $\rightarrow$ \code{cart}; \code{product_id} $\rightarrow$ \code{product} & \small CartItem\\
\small 13 & \small \code{orders} & \small \code{order_id} & \small \code{customer_id} $\rightarrow$ \code{customer} & \small Order\\
\small 14 & \small \code{order_item} & \small \code{order_id}, \code{product_id} & \small \code{order_id} $\rightarrow$ \code{orders}; \code{product_id} $\rightarrow$ \code{product} & \small OrderItem\\
\small 15 & \small \code{shipping} & \small \code{shipping_id} & \small \code{order_id} $\rightarrow$ \code{orders}; \code{address_id} $\rightarrow$ \code{address} & \small Shipping\\
\small 16 & \small \code{payment} & \small \code{payment_id} & \small \code{order_id} $\rightarrow$ \code{orders} & \small Payment\\
\small 17 & \small \code{pay_cash} & \small \code{payment_id} & \small \code{payment_id} $\rightarrow$ \code{payment} & \small PayCash\\
\small 18 & \small \code{pay_credit} & \small \code{payment_id} & \small \code{payment_id} $\rightarrow$ \code{payment} & \small PayCredit\\
\end{longtable}


\begin{longtable}{@{}P{5.4cm}P{2.3cm}P{2.0cm}P{5.3cm}@{}}
\caption{Foreign keys, relationship names and multiplicities}\label{tab:fks}\\
\toprule
\small \textbf{Foreign key} & \small \textbf{Name} & \small \textbf{Multiplicity} & \small \textbf{Meaning}\\
\midrule
\endfirsthead
\toprule
\small \textbf{Foreign key} & \small \textbf{Name} & \small \textbf{Multiplicity} & \small \textbf{Meaning}\\
\midrule
\endhead
\bottomrule
\endlastfoot
\small \code{customer_vip.customer_id} $\rightarrow$ \code{customer} & \small is a & \small 1 -- 0..1 & \small VIP part of a customer\\
\small \code{customer_new.customer_id} $\rightarrow$ \code{customer} & \small is a & \small 1 -- 0..1 & \small New-customer part of a customer\\
\small \code{address.customer_id} $\rightarrow$ \code{customer} & \small has & \small 1 -- 1..* & \small Addresses of a customer\\
\small \code{painting.product_id} $\rightarrow$ \code{product} & \small is a & \small 1 -- 0..1 & \small Painting part of a product\\
\small \code{painting.artist_id} $\rightarrow$ \code{artist} & \small paints & \small 1 -- 0..* & \small Artist of the painting\\
\small \code{frame.product_id} $\rightarrow$ \code{product} & \small is a & \small 1 -- 0..1 & \small Frame part of a product\\
\small \code{product_feature.product_id} $\rightarrow$ \code{product} & \small indexed by & \small 1 -- 0..1 & \small Feature vector used by search\\
\small \code{cart.customer_id} $\rightarrow$ \code{customer} & \small owns & \small 1 -- 0..1 & \small At most one cart per customer (UNIQUE)\\
\small \code{cart_item.cart_id} $\rightarrow$ \code{cart} & \small contains & \small 1 -- 0..* & \small Lines of a cart\\
\small \code{cart_item.product_id} $\rightarrow$ \code{product} & \small refers to & \small 1 -- 0..* & \small Product of a cart line\\
\small \code{orders.customer_id} $\rightarrow$ \code{customer} & \small places & \small 1 -- 0..* & \small Customer who placed the order\\
\small \code{order_item.order_id} $\rightarrow$ \code{orders} & \small contains & \small 1 -- 1..* & \small Lines of an order\\
\small \code{order_item.product_id} $\rightarrow$ \code{product} & \small refers to & \small 1 -- 0..* & \small Purchased product\\
\small \code{shipping.order_id} $\rightarrow$ \code{orders} & \small is delivered by & \small 1 -- 0..1 & \small At most one delivery per order (UNIQUE)\\
\small \code{shipping.address_id} $\rightarrow$ \code{address} & \small ships to & \small 1 -- 0..* & \small Delivery address\\
\small \code{payment.order_id} $\rightarrow$ \code{orders} & \small is paid by & \small 1 -- 0..* & \small Payment attempts of an order\\
\small \code{pay_cash.payment_id} $\rightarrow$ \code{payment} & \small is a & \small 1 -- 0..1 & \small Cash part of a payment\\
\small \code{pay_credit.payment_id} $\rightarrow$ \code{payment} & \small is a & \small 1 -- 0..1 & \small Card part of a payment\\
\end{longtable}


\section{Methods and directed relationships of the domain classes}\label{sec:B.3}

Table~\ref{tab:ops43} lists every method of the domain classes of Figure~\ref{fig:c43_design} with its full
signature and its responsibility, and the values of the three enumerations; Table~\ref{tab:nav43} gives the
direction and the role name of every association (Section~4.3).

\begin{longtable}{@{}P{2.6cm}P{7.0cm}P{5.4cm}@{}}
\caption{Methods of the domain classes at design level}\label{tab:ops43}\\
\toprule
\footnotesize \textbf{Class} & \footnotesize \textbf{Method} & \footnotesize \textbf{Responsibility}\\
\midrule
\endfirsthead
\toprule
\footnotesize \textbf{Class} & \footnotesize \textbf{Method} & \footnotesize \textbf{Responsibility}\\
\midrule
\endhead
\bottomrule
\endlastfoot
\footnotesize \code{FullName} & \footnotesize \code{+getDisplayName(): String} & \footnotesize Joins first, middle and last name for display\\
\footnotesize \code{Address} & \footnotesize \code{+format(): String} & \footnotesize Address on one line, printed on the order\\
\footnotesize \code{Customer} & \footnotesize \code{+getDisplayName(): String} & \footnotesize Display name taken from FullName\\
\footnotesize \code{Customer} & \footnotesize \code{+getDiscountRate(): BigDecimal} & \footnotesize Default discount 0, overridden by subclasses\\
\footnotesize \code{Customer} & \footnotesize \code{+addAddress(address: Address): void} & \footnotesize Adds a delivery address\\
\footnotesize \code{Customer} & \footnotesize \code{+findAddress(addressId: String): Address} & \footnotesize Finds an address by id; fails when it is not the customer's\\
\footnotesize \code{CustomerVIP} & \footnotesize \code{+getDiscountRate(): BigDecimal} & \footnotesize Discount rate of the VIP tier (polymorphism)\\
\footnotesize \code{CustomerVIP} & \footnotesize \code{+addPoints(points: int): void} & \footnotesize Adds reward points, refuses negative values\\
\footnotesize \code{CustomerNew} & \footnotesize \code{+hasWelcomeVoucher(): boolean} & \footnotesize Tells whether the new customer still has a voucher\\
\footnotesize \code{DetailKind}\newline\emph{«enumeration»} & \footnotesize \code{PAINTING}, \code{FRAME} & \footnotesize Valid status values\\
\footnotesize \code{Category} & \footnotesize \code{+accepts(product: Product): boolean} & \footnotesize Tells whether a product has the kind of detail the category names\\
\footnotesize \code{Product} & \footnotesize \code{+isAvailable(quantity: int): boolean} & \footnotesize Checks that enough stock is left for a quantity\\
\footnotesize \code{Product} & \footnotesize \code{+decreaseStock(quantity: int): void} & \footnotesize Decreases stock when ordered; refuses to oversell\\
\footnotesize \code{Product} & \footnotesize \code{+getDetailKind(): DetailKind} \{abstract\} & \footnotesize Kind of detail of the product, decided by the subclass (abstract)\\
\footnotesize \code{Product} & \footnotesize \code{+keywords(): String} & \footnotesize Keywords used by text search\\
\footnotesize \code{Artist} & \footnotesize \code{+getDisplayName(): String} & \footnotesize Artist name with years of birth and death\\
\footnotesize \code{Painting} & \footnotesize \code{+getDetailKind(): DetailKind} & \footnotesize Returns PAINTING\\
\footnotesize \code{Painting} & \footnotesize \code{+keywords(): String} & \footnotesize Adds style, subject, colour, medium and artist name to the keywords\\
\footnotesize \code{Painting} & \footnotesize \code{+sizeLabel(): String} & \footnotesize Print size as "width x height cm"\\
\footnotesize \code{Frame} & \footnotesize \code{+getDetailKind(): DetailKind} & \footnotesize Returns FRAME\\
\footnotesize \code{Frame} & \footnotesize \code{+keywords(): String} & \footnotesize Adds frame material and colour to the keywords\\
\footnotesize \code{Cart} & \footnotesize \code{+addItem(product: Product, quantity: int): void} & \footnotesize Adds a product; merges with an existing line of the same product\\
\footnotesize \code{Cart} & \footnotesize \code{+updateQuantity(productId: String, quantity: int): void} & \footnotesize Changes a quantity; 0 removes the line\\
\footnotesize \code{Cart} & \footnotesize \code{+removeItem(productId: String): void} & \footnotesize Removes one line\\
\footnotesize \code{Cart} & \footnotesize \code{+calculateTotal(): BigDecimal} & \footnotesize Provisional total of the cart\\
\footnotesize \code{Cart} & \footnotesize \code{+isEmpty(): boolean} & \footnotesize Whether the cart is empty (checkout refuses an empty cart)\\
\footnotesize \code{Cart} & \footnotesize \code{+clear(): void} & \footnotesize Empties the cart after a successful checkout\\
\footnotesize \code{CartItem} & \footnotesize \code{+CartItem(product: Product, quantity: int)} & \footnotesize Creates a line and records the price when added\\
\footnotesize \code{CartItem} & \footnotesize \code{+getSubtotal(): BigDecimal} & \footnotesize Amount of the cart line\\
\footnotesize \code{CartItem} & \footnotesize \code{+increase(quantity: int): void} & \footnotesize Increases the quantity when the same product is added again\\
\footnotesize \code{OrderStatus}\newline\emph{«enumeration»} & \footnotesize \code{PLACED}, \code{PROCESSING}, \code{SHIPPED}, \code{DELIVERED}, \code{CANCELLED} & \footnotesize Valid status values\\
\footnotesize \code{Order} & \footnotesize \code{+addItem(item: OrderItem): void} & \footnotesize Adds an order line\\
\footnotesize \code{Order} & \footnotesize \code{+addPayment(payment: Payment): void} & \footnotesize Attaches a payment attempt\\
\footnotesize \code{Order} & \footnotesize \code{+calculateTotal(discountRate: BigDecimal): BigDecimal} & \footnotesize Total = subtotal - discount + shipping fee, never negative\\
\footnotesize \code{Order} & \footnotesize \code{+cancel(): void} & \footnotesize Cancels the order; allowed only when PLACED or PROCESSING\\
\footnotesize \code{Order} & \footnotesize \code{+advanceTo(next: OrderStatus): void} & \footnotesize Moves the status in the order PLACED, PROCESSING, SHIPPED, DELIVERED\\
\footnotesize \code{Order} & \footnotesize \code{+isOwnedBy(customerId: String): boolean} & \footnotesize Checks that the order belongs to the signed-in customer\\
\footnotesize \code{OrderItem} & \footnotesize \code{+fromCartItem(item: CartItem, current: Product): OrderItem} \{static\} & \footnotesize Turns a cart line into an order line at the current catalogue price\\
\footnotesize \code{OrderItem} & \footnotesize \code{+getSubtotal(): BigDecimal} & \footnotesize Amount of the order line at the frozen price\\
\footnotesize \code{ShippingStatus}\newline\emph{«enumeration»} & \footnotesize \code{PREPARING}, \code{IN_TRANSIT}, \code{DELIVERED} & \footnotesize Valid status values\\
\footnotesize \code{Shipping} & \footnotesize \code{+updateStatus(newStatus: ShippingStatus): void} & \footnotesize Changes the delivery status\\
\footnotesize \code{Shipping} & \footnotesize \code{+markDelivered(date: LocalDate): void} & \footnotesize Records the delivery date and DELIVERED\\
\footnotesize \code{PaymentStatus}\newline\emph{«enumeration»} & \footnotesize \code{PENDING}, \code{SUCCEEDED}, \code{FAILED} & \footnotesize Valid status values\\
\footnotesize \code{Payment} & \footnotesize \code{+process(): boolean} \{abstract\} & \footnotesize Processes the payment, each method in its own way (abstract)\\
\footnotesize \code{Payment} & \footnotesize \code{+markSucceeded(): void} & \footnotesize Marks success and records the time\\
\footnotesize \code{Payment} & \footnotesize \code{+markFailed(): void} & \footnotesize Marks failure\\
\footnotesize \code{PayCash} & \footnotesize \code{+process(): boolean} & \footnotesize Checks the cash received and computes the change\\
\footnotesize \code{PayCredit} & \footnotesize \code{+process(): boolean} & \footnotesize Succeeds only with an authorisation code from the payment service\\
\footnotesize \code{PayCredit} & \footnotesize \code{+maskedCard(): String} & \footnotesize Shows the card as **** **** **** 1234\\
\end{longtable}


\begin{longtable}{@{}P{3.5cm}P{3.8cm}P{2.0cm}P{5.7cm}@{}}
\caption{Directed relationships at design level}\label{tab:nav43}\\
\toprule
\footnotesize \textbf{Relationship} & \footnotesize \textbf{Direction and role} & \footnotesize \textbf{Multiplicity} & \footnotesize \textbf{Meaning}\\
\midrule
\endfirsthead
\toprule
\footnotesize \textbf{Relationship} & \footnotesize \textbf{Direction and role} & \footnotesize \textbf{Multiplicity} & \footnotesize \textbf{Meaning}\\
\midrule
\endhead
\bottomrule
\endlastfoot
\footnotesize \code{Customer} -- \code{FullName} (has) & \footnotesize \code{Customer -> fullName} & \footnotesize 1 -- 1 & \footnotesize Each customer has exactly one name; the name does not exist without the customer, hence composition.\\
\footnotesize \code{Customer} -- \code{Address} (has) & \footnotesize \code{Customer -> addresses} & \footnotesize 1 -- 1..* & \footnotesize A customer keeps at least one delivery address; Shipping also uses it, so only aggregation.\\
\footnotesize \code{Cart} -- \code{Customer} (owned by) & \footnotesize \code{Cart -> owner} & \footnotesize 0..1 -- 1 & \footnotesize At most one active cart per customer; a new customer may have none.\\
\footnotesize \code{Cart} -- \code{CartItem} (contains) & \footnotesize \code{Cart -> items} & \footnotesize 1 -- 0..* & \footnotesize A cart line exists only inside its cart; an empty cart is valid.\\
\footnotesize \code{CartItem} -- \code{Product} (refers to) & \footnotesize \code{CartItem -> product} & \footnotesize 0..* -- 1 & \footnotesize Each cart line refers to exactly one product; a product appears in many carts.\\
\footnotesize \code{Order} -- \code{Customer} (placed by) & \footnotesize \code{Order -> customer} & \footnotesize 0..* -- 1 & \footnotesize A customer places many orders; a new customer has none yet.\\
\footnotesize \code{Order} -- \code{OrderItem} (contains) & \footnotesize \code{Order -> items} & \footnotesize 1 -- 1..* & \footnotesize An order line is part of its order; an order without lines makes no sense.\\
\footnotesize \code{OrderItem} -- \code{Product} (refers to) & \footnotesize \code{OrderItem -> product} & \footnotesize 0..* -- 1 & \footnotesize Each order line refers to exactly one product; a product appears in many orders.\\
\footnotesize \code{Order} -- \code{Shipping} (delivered by) & \footnotesize \code{Order -> shipping} & \footnotesize 1 -- 0..1 & \footnotesize A new order may not be shipped yet; once shipped it has exactly one delivery record.\\
\footnotesize \code{Shipping} -- \code{Address} (ships to) & \footnotesize \code{Shipping -> address} & \footnotesize 0..* -- 1 & \footnotesize Each delivery goes to exactly one address; an address receives many deliveries.\\
\footnotesize \code{Order} -- \code{Payment} (paid by) & \footnotesize \code{Order -> payments} & \footnotesize 1 -- 0..* & \footnotesize An order may be unpaid, or have several attempts when one fails.\\
\footnotesize \code{Painting} -- \code{Artist} (painted by) & \footnotesize \code{Painting -> artist} & \footnotesize 0..* -- 1 & \footnotesize An artist has many paintings in the gallery; each print belongs to exactly one artist.\\
\footnotesize \code{Product} -- \code{Category} (classified in) & \footnotesize \code{Product -> category} & \footnotesize 0..* -- 1 & \footnotesize \\
\end{longtable}


\section{Layers, packages and classes}\label{sec:B.4}

Table~\ref{tab:layers} lists the packages of each layer of Figure~\ref{fig:c44_layers} with their classes
(Section~4.4).

\begin{longtable}{@{}P{2.1cm}P{4.4cm}P{4.8cm}P{3.7cm}@{}}
\caption{Layers, packages and classes of Họa.}\label{tab:layers}\\
\toprule
\footnotesize \textbf{Layer} & \footnotesize \textbf{Package} & \footnotesize \textbf{Classes} & \footnotesize \textbf{Responsibility}\\
\midrule
\endfirsthead
\toprule
\footnotesize \textbf{Layer} & \footnotesize \textbf{Package} & \footnotesize \textbf{Classes} & \footnotesize \textbf{Responsibility}\\
\midrule
\endhead
\bottomrule
\endlastfoot
\footnotesize Presentation & \footnotesize \code{com.ptit.ecom.presentation.view} & \footnotesize \code{SearchUI}, \code{VoiceInput}, \code{RoomPreview}, \code{ImageUpload}, \code{SearchResultView}, \code{CartView}, \code{OrderView}, \code{AssistantView}, \code{AdminView} & \footnotesize Browser components (HTML, CSS, JS) named after the A06 presentation layer\\
\footnotesize Presentation & \footnotesize \code{com.ptit.ecom.presentation.controller} & \footnotesize \code{SearchController}, \code{CartController}, \code{CheckoutController}, \code{OrderController}, \code{AdminController}, \code{AssistantController} & \footnotesize Receive decoded requests, call services, return results\\
\footnotesize Presentation & \footnotesize \code{com.ptit.ecom.presentation.web} & \footnotesize \code{SessionManager}, \code{WebServer}, \code{EcomApplication} & \footnotesize HTTP server, sessions, composition root\\
\footnotesize Application & \footnotesize \code{com.ptit.ecom.application.service} & \footnotesize \code{CartService}, \code{CheckoutService}, \code{OrderService}, \code{ProductService}, \code{AssistantReply}, \code{AssistantService} & \footnotesize Coordinate the cart, checkout, order, product and assistant use cases\\
\footnotesize Application & \footnotesize \code{com.ptit.ecom.application.search} & \footnotesize \code{CommonQuery}, \code{SearchResult}, \code{QueryService}, \code{SpeechService}, \code{ImageService}, \code{RankingService}, \code{SearchService} & \footnotesize Multimodal search: common query, speech, image, ranking\\
\footnotesize Data & \footnotesize \code{com.ptit.ecom.infrastructure.repository} & \footnotesize \code{CustomerRepository}, \code{ProductRepository}, \code{CartRepository}, \code{OrderRepository} & \footnotesize Repository interfaces used by the application layer\\
\footnotesize Data & \footnotesize \code{com.ptit.ecom.infrastructure.repository.sqlite} & \footnotesize \code{Database}, \code{SqliteCustomerRepository}, \code{SqliteProductRepository}, \code{SqliteCartRepository}, \code{SqliteOrderRepository} & \footnotesize JDBC implementations of the repositories on a SQLite file\\
\footnotesize Data & \footnotesize \code{com.ptit.ecom.infrastructure.search} & \footnotesize \code{VectorIndex} & \footnotesize Vector index and cosine similarity\\
\footnotesize Data & \footnotesize \code{com.ptit.ecom.infrastructure.ai} & \footnotesize \code{ToolCall}, \code{LlmReply}, \code{LlmClient}, \code{SpeechClient}, \code{WhisperClient}, \code{GemmaClient} & \footnotesize Clients of Gemma 4 (OpenAI-compatible) and of the Whisper voice server\\
\footnotesize Domain model & \footnotesize \code{com.ptit.ecom.domain.customer} & \footnotesize \code{FullName}, \code{Address}, \code{Customer}, \code{CustomerVIP}, \code{CustomerNew} & \footnotesize Customer, name, address, customer tiers\\
\footnotesize Domain model & \footnotesize \code{com.ptit.ecom.domain.product} & \footnotesize \code{Product}, \code{Artist}, \code{Painting}, \code{Frame} & \footnotesize Product, painting, frame, artist\\
\footnotesize Domain model & \footnotesize \code{com.ptit.ecom.domain.cart} & \footnotesize \code{Cart}, \code{CartItem} & \footnotesize Cart and cart lines\\
\footnotesize Domain model & \footnotesize \code{com.ptit.ecom.domain.order} & \footnotesize \code{OrderStatus}, \code{Order}, \code{OrderItem}, \code{ShippingStatus}, \code{Shipping} & \footnotesize Order, order lines, shipping and their statuses\\
\footnotesize Domain model & \footnotesize \code{com.ptit.ecom.domain.payment} & \footnotesize \code{PaymentStatus}, \code{Payment}, \code{PayCash}, \code{PayCredit} & \footnotesize Payments and payment status\\
\end{longtable}


\section{Traceability from requirements to design, data and code}\label{sec:B.5}

Table~\ref{tab:trace46} follows the main requirements of Chapter~\ref{ch:req} through the artefacts of Chapter~4 --
analysis classes, tables, design classes and methods, packages -- to the code (Section~4.6).

\begin{longtable}{@{}P{1.4cm}P{3.0cm}P{2.6cm}P{4.6cm}P{3.4cm}@{}}
\caption{Traceability from requirements to design, data and code}\label{tab:trace46}\\
\toprule
\textbf{Req.} & \textbf{Analysis classes} & \textbf{Tables} & \textbf{Design classes and methods} & \textbf{Package}\\
\midrule
\endfirsthead
\toprule
\textbf{Req.} & \textbf{Analysis classes} & \textbf{Tables} & \textbf{Design classes and methods} & \textbf{Package}\\
\midrule
\endhead
\small FR01 & \small \code{Product}, \code{Painting}, \code{Artist} & \small \code{product}, \code{painting}, \code{artist} & \small \code{QueryService.textQuery}, \code{SearchService.searchText}, \code{RankingService.rank} & \small \code{application.search}\\
\small FR03 & \small \code{Painting} & \small \code{painting}, \code{product_feature} & \small \code{ImageService.embed}, \code{VectorIndex.cosineSimilarity} & \small \code{application.search}, \code{infrastructure.search}\\
\small FR06 & \small \code{Cart}, \code{CartItem} & \small \code{cart}, \code{cart_item} & \small \code{Cart.addItem}, \code{CartService.addToCart} & \small \code{domain.cart}, \code{application.service}\\
\small FR07 & \small \code{Cart}, \code{Order}, \code{OrderItem}, \code{Shipping}, \code{Payment} & \small \code{orders}, \code{order_item}, \code{shipping}, \code{payment} & \small \code{CheckoutService.checkout}, \code{OrderItem.fromCartItem}, \code{Order.calculateTotal} & \small \code{application.service}, \code{domain.order}\\
\small FR09 & \small \code{Order}, \code{Customer} & \small \code{orders.customer_id} & \small \code{OrderService.getOrder}, \code{Order.isOwnedBy}, \code{OrderRepository.findOrder} & \small \code{application.service}, \code{infrastructure.repository}\\
\small FR11 & \small \code{Order} & \small \code{orders.status} & \small \code{Order.cancel}, \code{OrderService.cancelOrder} & \small \code{domain.order}\\
\small FR12 & \small \code{Product}, \code{Painting} & \small \code{product}, \code{painting} & \small \code{ProductService.updatePriceAndStock}, \code{ProductService.addPainting} & \small \code{application.service}\\
\small FR13 & \small \code{Order} & \small \code{orders.status} & \small \code{Order.advanceTo}, \code{OrderService.updateStatus} & \small \code{domain.order}\\
\small FR14 & \small \code{Painting} & \small \code{painting} & \small \code{SearchService.describeRoom}, \code{SearchService.searchRoom} & \small \code{application.search}\\
\small NFR05 & \small \code{PayCredit} & \small \code{pay_credit.card_last4} & \small \code{PayCredit.maskedCard} & \small \code{domain.payment}\\
\bottomrule
\end{longtable}

\section{Mapping from the UML model to Java}\label{sec:B.6}

Table~\ref{tab:uml2java} lists how each element of the design model of Sections~4.3 and~4.4 becomes Java code in the
skeleton generated by \code{tools/gen_java.py} (Section~4.5).

\begin{longtable}{@{}P{4.4cm}P{5.3cm}P{5.3cm}@{}}
\caption{Mapping from the UML model to Java}\label{tab:uml2java}\\
\toprule
\textbf{UML element (4.3, 4.4)} & \textbf{Java} & \textbf{Example}\\
\midrule
\endfirsthead
\toprule
\textbf{UML element (4.3, 4.4)} & \textbf{Java} & \textbf{Example}\\
\midrule
\endhead
\small Layer package and sub-package & \small Package \code{com.ptit.ecom.<layer>.<area>} & \small \code{com.ptit.ecom.domain.order}\\
\small Class, abstract class & \small \code{public class}, \code{public abstract class} & \small \code{Product}, \code{Payment}\\
\small Interface & \small \code{public interface} & \small \code{OrderRepository}, \code{LlmClient}\\
\small Enumeration & \small \code{public enum} & \small \code{OrderStatus}\\
\small Private attribute & \small Private field with the same name and type & \small \code{-totalAmount: BigDecimal}\\
\small Association end with role & \small Field named after the role; \code{List} for an upper bound * & \small \code{Order.items}, \code{Painting.artist}\\
\small Generalization & \small \code{extends} & \small \code{Painting extends Product}\\
\small Realization & \small \code{implements} & \small \code{SqliteOrderRepository implements OrderRepository}\\
\small «use» dependency & \small Field set by the constructor & \small \code{cartRepository} of \code{CheckoutService}\\
\small Operation & \small Method with the same name, parameters and return type & \small \code{OrderService.getOrder}\\
\bottomrule
\end{longtable}

\section{Report of the code-model check}\label{sec:B.7}

The report of \code{tools/check_code_model.py} (the second part of \code{build/test_results.txt}) compares the Java
classes, the JavaScript views, the real database and the SQL of the repositories with \code{tools/model.py}. The runs
of per-class ``OK'' lines are shortened to one counted line; every other line is reproduced unchanged.

\txtfileor{generated/model_check.txt}{../build/test_results.txt}

\section{Mapping of design concepts to technology}\label{sec:B.8}

Table~\ref{tab:techmap} records the technology chosen for each design concept of Section~4.4 (Lecture~05).

\begin{longtable}{@{}P{3.2cm}P{5.6cm}P{6.2cm}@{}}
\caption{Mapping of design concepts to technology}\label{tab:techmap}\\
\toprule
\textbf{Design concept} & \textbf{Technology} & \textbf{Where}\\
\midrule
\endfirsthead
\toprule
\textbf{Design concept} & \textbf{Technology} & \textbf{Where}\\
\midrule
\endhead
\small View & \small HTML, CSS and one JavaScript module per view class & \small \code{demo/web}\\
\small Web server, session & \small Java HTTP server of the JDK, JSON & \small \code{WebServer}, \code{SessionManager}\\
\small Controller & \small Plain Java class called by \code{WebServer} & \small \code{presentation.controller}\\
\small Service & \small Plain Java class & \small \code{application.service}, \code{application.search}\\
\small Domain entity & \small Plain Java object (no framework annotations) & \small \code{domain.*}\\
\small Repository & \small Java interface + JDBC implementation & \small \code{infrastructure.repository}, \code{.sqlite}\\
\small Database & \small SQLite file & \small \code{demo/data/ecom.db}\\
\small Language model & \small Gemma~4 through the OpenAI-compatible HTTP API & \small \code{GemmaClient}\\
\small Speech-to-text & \small Self-hosted faster-whisper server (Python) & \small \code{WhisperClient}, \code{demo/voice-server}\\
\bottomrule
\end{longtable}

\chapter{Requirement Reference Tables}\label{app:reqtables}\label{app:C}

The tables of this appendix complete Chapter~1: how each non-functional requirement is verified (Section~1.9), the
full traceability matrix between requirements and use cases (Section~1.12) and the relationships of the use case
diagram (Section~1.10). Both are generated from
\code{tools/model.py}.

\section{Verification of the non-functional requirements}\label{sec:C.1}

\begin{longtable}{@{}P{1.3cm}P{5.0cm}P{8.7cm}@{}}
\caption{How the non-functional requirements are verified}\label{tab:nfrcheck}\\
\toprule
\small \textbf{ID} & \small \textbf{Short statement} & \small \textbf{Verification}\\
\midrule
\endfirsthead
\toprule
\small \textbf{ID} & \small \textbf{Short statement} & \small \textbf{Verification}\\
\midrule
\endhead
\bottomrule
\endlastfoot
\small \textbf{NFR01} & \small Security: the customer id comes from the session and order ownership is checked on the server & \small Customer C002 asks for order 20261001 of C001 and is refused; the customer id is never read from the request body.\\
\small \textbf{NFR02} & \small Fallback: when Gemma fails or times out, search falls back to keyword rules & \small With a wrong Gemma address the search still returns results from keyword rules and the page says that the fallback is used.\\
\small \textbf{NFR03} & \small Uploaded images: JPEG, PNG or WebP only, at most 5 MB & \small A file that is not an image, or is larger than 5 MB, is rejected before Gemma is called.\\
\small \textbf{NFR04} & \small Layered architecture: the UI never accesses the database, AI sits behind the LlmClient interface & \small Views call controllers only; repositories, GemmaClient and WhisperClient are created only in the composition root.\\
\small \textbf{NFR05} & \small Payment data: never store the full card number, only the last four digits and the authorisation code & \small Table pay\_credit has card\_last4 and authorization\_code and no column for the full card number.\\
\small \textbf{NFR06} & \small Runs offline on one machine: SQLite database, network only for the Gemma server & \small The application runs on one machine with a SQLite file; its only outside connection is the Gemma 4 server.\\
\small \textbf{NFR07} & \small Performance: text search answers within 2 seconds on the demo data; an AI call waits at most 15 seconds & \small searchText is timed on the demo data; a Gemma call is cut after the timeoutSeconds of GemmaClient.\\
\small \textbf{NFR08} & \small Voice privacy: audio is processed only on the machine by self-hosted Whisper; a 5-second sentence is transcribed within 3 seconds & \small Audio is sent only to the Whisper server on 127.0.0.1; the time to transcribe a 5-second sentence is measured (build/eval/voice\_eval.txt).\\
\small \textbf{NFR09} & \small Usability: every search, cart and order action ends in a visible result, empty state or error that says what to do & \small tools/e2e\_check.py checks the answers of an empty cart, a foreign address, another customer's order and a shipped order; the browser shows them as messages (Chapter 5 screenshots).\\
\small \textbf{NFR10} & \small Maintainability: class, method, table and column names of the code match the design model, checked automatically & \small tools/check\_code\_model.py compares every class, method, parameter, table, column and SQL statement of the code with tools/model.py; demo/test.sh runs it.\\
\small \textbf{NFR11} & \small Data integrity: keys and constraints enforced, stock never negative, a product's details in the table its category names & \small tools/gen\_sql.py runs statements that break a key, a CHECK or a category and expects SQLite to refuse them (build/sql\_validation.txt); stock is decreased by a conditional UPDATE; tests refuse a product its category does not accept.\\
\small \textbf{NFR12} & \small Availability: catalogue, cart, checkout, orders and administration work without the AI services & \small tools/e2e\_check.py starts a second server whose Gemma address is unreachable and places an order with it (build/e2e\_results.txt).\\
\end{longtable}


\section{Requirement -- use case traceability matrix}\label{sec:C.2}

{\setlength{\tabcolsep}{2.2pt}
\begin{longtable}{@{}P{1.2cm}>{\centering\arraybackslash}p{0.74cm}>{\centering\arraybackslash}p{0.74cm}>{\centering\arraybackslash}p{0.74cm}>{\centering\arraybackslash}p{0.74cm}>{\centering\arraybackslash}p{0.74cm}>{\centering\arraybackslash}p{0.74cm}>{\centering\arraybackslash}p{0.74cm}>{\centering\arraybackslash}p{0.74cm}>{\centering\arraybackslash}p{0.74cm}>{\centering\arraybackslash}p{0.74cm}>{\centering\arraybackslash}p{0.74cm}>{\centering\arraybackslash}p{0.74cm}>{\centering\arraybackslash}p{0.74cm}>{\centering\arraybackslash}p{0.74cm}@{}}
\caption{Requirement -- use case traceability matrix}\label{tab:traceuc}\\
\toprule
\scriptsize \textbf{ID} & \scriptsize \textbf{\scriptsize UC01} & \scriptsize \textbf{\scriptsize UC02} & \scriptsize \textbf{\scriptsize UC03} & \scriptsize \textbf{\scriptsize UC04} & \scriptsize \textbf{\scriptsize UC05} & \scriptsize \textbf{\scriptsize UC06} & \scriptsize \textbf{\scriptsize UC07} & \scriptsize \textbf{\scriptsize UC08} & \scriptsize \textbf{\scriptsize UC09} & \scriptsize \textbf{\scriptsize UC10} & \scriptsize \textbf{\scriptsize UC11} & \scriptsize \textbf{\scriptsize UC12} & \scriptsize \textbf{\scriptsize UC13} & \scriptsize \textbf{\scriptsize UC14}\\
\midrule
\endfirsthead
\toprule
\scriptsize \textbf{ID} & \scriptsize \textbf{\scriptsize UC01} & \scriptsize \textbf{\scriptsize UC02} & \scriptsize \textbf{\scriptsize UC03} & \scriptsize \textbf{\scriptsize UC04} & \scriptsize \textbf{\scriptsize UC05} & \scriptsize \textbf{\scriptsize UC06} & \scriptsize \textbf{\scriptsize UC07} & \scriptsize \textbf{\scriptsize UC08} & \scriptsize \textbf{\scriptsize UC09} & \scriptsize \textbf{\scriptsize UC10} & \scriptsize \textbf{\scriptsize UC11} & \scriptsize \textbf{\scriptsize UC12} & \scriptsize \textbf{\scriptsize UC13} & \scriptsize \textbf{\scriptsize UC14}\\
\midrule
\endhead
\bottomrule
\endlastfoot
\scriptsize \textbf{FR01} & \scriptsize  & \scriptsize $\bullet$ & \scriptsize  & \scriptsize  & \scriptsize  & \scriptsize  & \scriptsize  & \scriptsize  & \scriptsize  & \scriptsize  & \scriptsize  & \scriptsize  & \scriptsize  & \scriptsize \\
\scriptsize \textbf{FR02} & \scriptsize  & \scriptsize  & \scriptsize $\bullet$ & \scriptsize  & \scriptsize  & \scriptsize  & \scriptsize  & \scriptsize  & \scriptsize  & \scriptsize  & \scriptsize  & \scriptsize  & \scriptsize  & \scriptsize \\
\scriptsize \textbf{FR03} & \scriptsize  & \scriptsize  & \scriptsize  & \scriptsize $\bullet$ & \scriptsize  & \scriptsize  & \scriptsize  & \scriptsize  & \scriptsize  & \scriptsize  & \scriptsize  & \scriptsize  & \scriptsize  & \scriptsize \\
\scriptsize \textbf{FR04} & \scriptsize $\bullet$ & \scriptsize  & \scriptsize  & \scriptsize  & \scriptsize  & \scriptsize  & \scriptsize  & \scriptsize  & \scriptsize  & \scriptsize  & \scriptsize  & \scriptsize  & \scriptsize  & \scriptsize \\
\scriptsize \textbf{FR05} & \scriptsize  & \scriptsize  & \scriptsize  & \scriptsize  & \scriptsize $\bullet$ & \scriptsize  & \scriptsize  & \scriptsize  & \scriptsize  & \scriptsize  & \scriptsize  & \scriptsize  & \scriptsize  & \scriptsize \\
\scriptsize \textbf{FR06} & \scriptsize  & \scriptsize  & \scriptsize  & \scriptsize  & \scriptsize  & \scriptsize $\bullet$ & \scriptsize  & \scriptsize  & \scriptsize  & \scriptsize  & \scriptsize  & \scriptsize  & \scriptsize  & \scriptsize \\
\scriptsize \textbf{FR07} & \scriptsize  & \scriptsize  & \scriptsize  & \scriptsize  & \scriptsize  & \scriptsize  & \scriptsize $\bullet$ & \scriptsize  & \scriptsize  & \scriptsize  & \scriptsize  & \scriptsize  & \scriptsize  & \scriptsize \\
\scriptsize \textbf{FR08} & \scriptsize  & \scriptsize  & \scriptsize  & \scriptsize  & \scriptsize  & \scriptsize  & \scriptsize  & \scriptsize $\bullet$ & \scriptsize  & \scriptsize  & \scriptsize  & \scriptsize  & \scriptsize  & \scriptsize \\
\scriptsize \textbf{FR09} & \scriptsize  & \scriptsize  & \scriptsize  & \scriptsize  & \scriptsize  & \scriptsize  & \scriptsize  & \scriptsize  & \scriptsize $\bullet$ & \scriptsize  & \scriptsize  & \scriptsize  & \scriptsize  & \scriptsize \\
\scriptsize \textbf{FR10} & \scriptsize  & \scriptsize  & \scriptsize  & \scriptsize  & \scriptsize  & \scriptsize  & \scriptsize  & \scriptsize  & \scriptsize  & \scriptsize $\bullet$ & \scriptsize  & \scriptsize  & \scriptsize  & \scriptsize \\
\scriptsize \textbf{FR11} & \scriptsize  & \scriptsize  & \scriptsize  & \scriptsize  & \scriptsize  & \scriptsize  & \scriptsize  & \scriptsize  & \scriptsize  & \scriptsize  & \scriptsize $\bullet$ & \scriptsize  & \scriptsize  & \scriptsize \\
\scriptsize \textbf{FR12} & \scriptsize  & \scriptsize  & \scriptsize  & \scriptsize  & \scriptsize  & \scriptsize  & \scriptsize  & \scriptsize  & \scriptsize  & \scriptsize  & \scriptsize  & \scriptsize $\bullet$ & \scriptsize  & \scriptsize \\
\scriptsize \textbf{FR13} & \scriptsize  & \scriptsize  & \scriptsize  & \scriptsize  & \scriptsize  & \scriptsize  & \scriptsize  & \scriptsize  & \scriptsize  & \scriptsize  & \scriptsize  & \scriptsize  & \scriptsize $\bullet$ & \scriptsize \\
\scriptsize \textbf{FR14} & \scriptsize  & \scriptsize  & \scriptsize  & \scriptsize  & \scriptsize  & \scriptsize  & \scriptsize  & \scriptsize  & \scriptsize  & \scriptsize  & \scriptsize  & \scriptsize  & \scriptsize  & \scriptsize $\bullet$\\
\scriptsize \textbf{NFR01} & \scriptsize  & \scriptsize  & \scriptsize  & \scriptsize  & \scriptsize  & \scriptsize  & \scriptsize  & \scriptsize $\bullet$ & \scriptsize $\bullet$ & \scriptsize $\bullet$ & \scriptsize  & \scriptsize  & \scriptsize  & \scriptsize \\
\scriptsize \textbf{NFR02} & \scriptsize  & \scriptsize $\bullet$ & \scriptsize $\bullet$ & \scriptsize $\bullet$ & \scriptsize  & \scriptsize  & \scriptsize  & \scriptsize  & \scriptsize  & \scriptsize $\bullet$ & \scriptsize  & \scriptsize  & \scriptsize  & \scriptsize \\
\scriptsize \textbf{NFR03} & \scriptsize  & \scriptsize  & \scriptsize  & \scriptsize $\bullet$ & \scriptsize  & \scriptsize  & \scriptsize  & \scriptsize  & \scriptsize  & \scriptsize  & \scriptsize  & \scriptsize  & \scriptsize  & \scriptsize \\
\scriptsize \textbf{NFR04} & \scriptsize $\bullet$ & \scriptsize  & \scriptsize  & \scriptsize  & \scriptsize  & \scriptsize  & \scriptsize  & \scriptsize  & \scriptsize  & \scriptsize  & \scriptsize  & \scriptsize  & \scriptsize  & \scriptsize \\
\scriptsize \textbf{NFR05} & \scriptsize  & \scriptsize  & \scriptsize  & \scriptsize  & \scriptsize  & \scriptsize  & \scriptsize $\bullet$ & \scriptsize  & \scriptsize  & \scriptsize  & \scriptsize  & \scriptsize  & \scriptsize  & \scriptsize \\
\scriptsize \textbf{NFR06} & \scriptsize $\bullet$ & \scriptsize  & \scriptsize  & \scriptsize  & \scriptsize  & \scriptsize  & \scriptsize  & \scriptsize  & \scriptsize  & \scriptsize  & \scriptsize  & \scriptsize  & \scriptsize  & \scriptsize \\
\scriptsize \textbf{NFR07} & \scriptsize  & \scriptsize $\bullet$ & \scriptsize  & \scriptsize  & \scriptsize  & \scriptsize  & \scriptsize  & \scriptsize  & \scriptsize  & \scriptsize  & \scriptsize  & \scriptsize  & \scriptsize  & \scriptsize \\
\scriptsize \textbf{NFR08} & \scriptsize  & \scriptsize  & \scriptsize $\bullet$ & \scriptsize  & \scriptsize  & \scriptsize  & \scriptsize  & \scriptsize  & \scriptsize  & \scriptsize  & \scriptsize  & \scriptsize  & \scriptsize  & \scriptsize \\
\end{longtable}
}


\section{Relationships between use cases}\label{sec:C.3}

\begin{longtable}{@{}P{5.4cm}P{2.6cm}P{7.0cm}@{}}
\caption{Relationships between use cases}\label{tab:uclinks}\\
\toprule
\small \textbf{Relationship} & \small \textbf{Kind} & \small \textbf{Meaning}\\
\midrule
\endfirsthead
\toprule
\small \textbf{Relationship} & \small \textbf{Kind} & \small \textbf{Meaning}\\
\midrule
\endhead
\bottomrule
\endlastfoot
\small UC02 Search by Keyword $\rightarrow$ UC01 Search Product & \small Generalization & \small Searching with keywords is one way to find paintings\\
\small UC03 Search by Voice $\rightarrow$ UC01 Search Product & \small Generalization & \small Searching by voice is one way to find paintings\\
\small UC04 Search by Image $\rightarrow$ UC01 Search Product & \small Generalization & \small Searching with a photo is one way to find paintings\\
\small UC07 Checkout $\rightarrow$ UC06 Manage Cart & \small «include» & \small Checkout always reads and checks the cart\\
\small UC10 Ask Shopping Assistant $\rightarrow$ UC01 Search Product & \small «include» & \small The assistant calls the painting search tool\\
\small UC10 Ask Shopping Assistant $\rightarrow$ UC09 Search Order & \small «include» & \small The assistant calls the order look-up tool\\
\small UC05 View Product $\rightarrow$ UC01 Search Product & \small «extend» & \small From the results the customer may open a painting\\
\small UC11 Cancel Order $\rightarrow$ UC08 View Order & \small «extend» & \small While viewing an order that is not shipped, the customer may cancel it\\
\small UC14 Find Painting for My Room $\rightarrow$ UC01 Search Product & \small Generalization & \small Finding a painting for a room is one way to find paintings; the input is a room photo\\
\end{longtable}


\chapter{Implementation Reference Tables}\label{app:impltables}\label{app:D}

The two tables of this appendix complete Chapter~5. They are generated from the code when the essay is built: the
HTTP endpoints from \code{WebServer.java} (Section~5.5) and the settings from \code{demo/config/app.properties}
(Section~5.2), whose API key is never printed.

\section{HTTP endpoints}\label{sec:D.1}

\begin{longtable}{@{}P{5.6cm}P{5.0cm}P{2.6cm}P{1.8cm}@{}}
\caption{HTTP endpoints of WebServer, the controller method each one calls and its use case}\label{tab:endpoints}\\
\toprule
\footnotesize \textbf{Method and path} & \footnotesize \textbf{Controller method} & \footnotesize \textbf{Identity} & \footnotesize \textbf{Use case}\\
\midrule
\endfirsthead
\toprule
\footnotesize \textbf{Method and path} & \footnotesize \textbf{Controller method} & \footnotesize \textbf{Identity} & \footnotesize \textbf{Use case}\\
\midrule
\endhead
\bottomrule
\endlastfoot
\footnotesize \code{POST /api/login} & \footnotesize \code{SessionManager.login} & \footnotesize -- & \footnotesize --\\
\footnotesize \code{GET /api/session} & \footnotesize \code{SessionManager.customerOf} & \footnotesize -- & \footnotesize --\\
\footnotesize \code{GET /api/products} & \footnotesize \code{SearchController.search} & \footnotesize -- & \footnotesize UC01\\
\footnotesize \code{GET /api/search} & \footnotesize \code{SearchController.search} & \footnotesize -- & \footnotesize UC02\\
\footnotesize \code{POST /api/voice/transcribe} & \footnotesize \code{SearchController.transcribe} & \footnotesize -- & \footnotesize UC03\\
\footnotesize \code{POST /api/search/voice} & \footnotesize \code{SearchController.searchVoice} & \footnotesize -- & \footnotesize UC03\\
\footnotesize \code{POST /api/search/image} & \footnotesize \code{SearchController.searchImage} & \footnotesize -- & \footnotesize UC04\\
\footnotesize \code{POST /api/search/multimodal} & \footnotesize \code{SearchController.searchMultimodal} & \footnotesize -- & \footnotesize UC01\\
\footnotesize \code{POST /api/room/describe} & \footnotesize \code{SearchController.describeRoom} & \footnotesize -- & \footnotesize UC14\\
\footnotesize \code{POST /api/room/search} & \footnotesize \code{SearchController.searchRoom} & \footnotesize -- & \footnotesize UC14\\
\footnotesize \code{POST /api/assistant} & \footnotesize \code{AssistantController.ask} & \footnotesize customer session & \footnotesize UC10\\
\footnotesize \code{GET /api/cart} & \footnotesize \code{CartController.viewCart} & \footnotesize customer session & \footnotesize UC06\\
\footnotesize \code{POST /api/cart/items} & \footnotesize \code{CartController.addItem} & \footnotesize customer session & \footnotesize UC06\\
\footnotesize \code{PUT /api/cart/items} & \footnotesize \code{CartController.updateQuantity} & \footnotesize customer session & \footnotesize UC06\\
\footnotesize \code{DELETE /api/cart/items} & \footnotesize \code{CartController.removeItem} & \footnotesize customer session & \footnotesize UC06\\
\footnotesize \code{POST /api/checkout} & \footnotesize \code{CartController.viewCart}, \code{CheckoutController.checkout} & \footnotesize customer session & \footnotesize UC07\\
\footnotesize \code{POST /api/addresses} & \footnotesize \code{CheckoutController.addAddress} & \footnotesize customer session & \footnotesize UC07\\
\footnotesize \code{GET /api/orders} & \footnotesize \code{OrderController.listOrders} & \footnotesize customer session & \footnotesize UC08\\
\footnotesize \code{GET /api/orders/<id>} & \footnotesize \code{OrderController.viewOrder} & \footnotesize customer session & \footnotesize UC08, UC09\\
\footnotesize \code{POST /api/orders/<id>/cancel} & \footnotesize \code{OrderController.cancelOrder}, \code{OrderController.viewOrder} & \footnotesize customer session & \footnotesize UC11\\
\footnotesize \code{GET /api/admin/products} & \footnotesize \code{AdminController.listProducts} & \footnotesize administrator session & \footnotesize UC12\\
\footnotesize \code{PUT /api/admin/products/<id>} & \footnotesize \code{AdminController.updateProduct} & \footnotesize administrator session & \footnotesize UC12\\
\footnotesize \code{POST /api/admin/paintings} & \footnotesize \code{AdminController.addPainting} & \footnotesize administrator session & \footnotesize UC12\\
\footnotesize \code{GET /api/admin/orders} & \footnotesize \code{AdminController.listAllOrders} & \footnotesize administrator session & \footnotesize UC13\\
\footnotesize \code{PUT /api/admin/orders/<id>} & \footnotesize \code{AdminController.updateOrderStatus} & \footnotesize administrator session & \footnotesize UC13\\
\end{longtable}


\section{Configuration}\label{sec:D.2}

\begin{longtable}{@{}P{4.4cm}P{6.6cm}P{4.0cm}@{}}
\caption{Settings in demo/config/app.properties and the class that reads them}\label{tab:config}\\
\toprule
\footnotesize \textbf{Key} & \footnotesize \textbf{Value} & \footnotesize \textbf{Read by}\\
\midrule
\endfirsthead
\toprule
\footnotesize \textbf{Key} & \footnotesize \textbf{Value} & \footnotesize \textbf{Read by}\\
\midrule
\endhead
\bottomrule
\endlastfoot
\footnotesize \code{server.port} & \footnotesize \code{8686} & \footnotesize \code{WebServer}\\
\footnotesize \code{db.path} & \footnotesize \code{data/ecom.db} & \footnotesize \code{Database}\\
\footnotesize \code{gemma.baseUrl} & \footnotesize \code{https://api.deepinfra.com/v1/openai/} & \footnotesize \code{GemmaClient}\\
\footnotesize \code{gemma.apiKey} & \footnotesize \emph{(read at start-up, not printed)} & \footnotesize \code{GemmaClient}\\
\footnotesize \code{gemma.model} & \footnotesize \code{google/gemma-4-31B-it-turbo} & \footnotesize \code{GemmaClient}\\
\footnotesize \code{gemma.timeoutSeconds} & \footnotesize \code{15} & \footnotesize \code{GemmaClient}\\
\footnotesize \code{whisper.baseUrl} & \footnotesize \code{http://127.0.0.1:9000} & \footnotesize \code{WhisperClient}\\
\footnotesize \code{whisper.model} & \footnotesize \code{small} & \footnotesize \code{WhisperClient}\\
\footnotesize \code{whisper.timeoutSeconds} & \footnotesize \code{15} & \footnotesize \code{WhisperClient}\\
\footnotesize \code{ranking.alpha} & \footnotesize \code{0.45} & \footnotesize \code{RankingService}\\
\footnotesize \code{ranking.beta} & \footnotesize \code{0.45} & \footnotesize \code{RankingService}\\
\footnotesize \code{ranking.gamma} & \footnotesize \code{0.10} & \footnotesize \code{RankingService}\\
\end{longtable}

